\Needspace{10\baselineskip}
\section{Conclusion}
\label{sec:conclusion}

EdgeRock shows that strong reasoning does not require a memory that grows with every token. Once the history a model carries is fixed, memory stops being a property of the input and becomes a property of the model: it can be sized once, at design time, and relied on for as long as the model reasons.

This changes what can be built around a reasoning model. A robot can run perception and reasoning in one continuous loop without its memory budget drifting as the task goes on. A dedicated chip can hold a model's entire history in on-chip memory, because that history no longer depends on context length. An agent's memory no longer grows with each tool call it makes. Serving systems can plan capacity by counting requests rather than by guessing how long each will run.

Several directions follow. Scaling RMA to larger backbones, beginning with additional EdgeRock sizes now in preparation, will test whether constant memory keeps pace with stronger reasoning. Applying the same conversion to other hybrid architectures would let existing models adopt constant memory without pretraining from scratch. Anchors could extend beyond text to streams of images, audio, and sensor data, where a growing history is most costly. And because each anchor records its source position, what a model chose to remember can be inspected directly.
